
%\needspace{4\baselineskip}
\color{uniblue}{\subsection{Защита от повышения напряжения (ЗПН)}}
\color{black}

\begin{enumerate}[label=\arabic{section}.\arabic{subsection}.\arabic*, labelsep=4pt, leftmargin=0pt, itemindent=57pt]

\item
ЗПН отслеживает повышение напряжения на секции шин и формирует команду на отключение выключателя. В устройстве защита от повышения напряжения представлена независимыми функциями (ступенями) «ЗПН 1» и «ЗПН 2».

\item
Пусковые условия формируются от ИО напряжения, реагирующих на превышение хотя бы одного из трех линейных напряжений выше уставки, отсутствии условий блокировки и выбора определенного режима работы уставки SGF2.

\item
При превышении ИО напряжения выше уставки активируются соответствующие выходные сигналы срабатывания ИО:
\begin{itemize}
    \item «ЗПН 1: ИО Uab»;
    \item «ЗПН 1: ИО Ubc»;
    \item «ЗПН 1: ИО Uca».
\end{itemize}

\item
Посредством уставки SGF2 предоставляется возможность настройки работы функции в одном из трех режимов:
\begin{itemize}
    \item по превышению одного линейного напряжения;
    \item по превышению двух линейных напряжений;
    \item по превышению трех линейных напряжений.
\end{itemize}

\item
Обязательным условием формирования выходных сигналов пуска и срабатывания функции является отсутствие активного сигнала блокировки от неисправности цепей напряжения ТН «БНН: Срабатывание».

\item
При превышении ИО напряжения ОП выше уставки активируется соответствующий выходной сигнал срабатывания ИО «ЗПН 1: ИО U2», «ЗПН 2: ИО U2».

\item
Уставкой SGF3 в функции предусмотрен контроль симметрии по напряжению обратной последовательности. В режиме работы по трем линейным напряжениям и увеличению напряжения ОП, пуск и срабатывание функции будет заблокировано.

\item
При наличии пусковых условий, срабатывание функции происходит после отсчета выдержки времени, таймера T1.

\item
Ввод функций в работу на этапе параметрирования устройства осуществляется своим программным переключателем SGF1 «Ввод\_функции». Информация о введенной ступени отображается выходным сигналом «ЗПН1: Ввод», «ЗПН2: Ввод».

\item
Функции могут быть выведены из работы оперативно путем активации сигнала «ОВ ЗПН 1» «ОВ ЗПН 2» соответственно,  что характеризуется активным состоянием сигнала «ЗПН 1: Оперативный вывод», «ЗПН 2: Оперативный вывод».

\item
При активном состоянии входного сигнала «ЗПН 1 на сигн», «ЗПН 2 на сигн» действие соответствующей функции переводится на сигнал.

\item
Алгоритм работы функции «ЗПН 1» выполнен в соответствии с рисунком \ref{zpn:img1}. Алгоритм работы функции «ЗПН 2» аналогичный. В таблице \ref{zpn:tbl1} приведены параметры, необходимые для настройки функций.

\vspace{1mm}
\begin{figure}[H]
\centering
\includegraphics[width=1\textwidth,height=1\textheight,keepaspectratio]{\fbpath/113. ЗПН/_latex/img/zpn.pdf}
\captionof{figure}{Функциональная схема логики ЗПН 1}\label{zpn:img1}
\end{figure}

\small
\begin{longtable}{|>{\centering\arraybackslash}m{5.3cm}|>{\centering\arraybackslash}m{3.3cm}|>{\centering\arraybackslash}m{4.2cm}|>{\centering\arraybackslash}m{1.8cm}|>{\centering\arraybackslash}m{1cm}|}
\caption{Параметры для настройки ЗПН\hfill\vspace{-0.5\baselineskip}}\label{zpn:tbl1}\\ 
\hline
\rowcolor{gray!30}
Параметр (Параметр на ИЧМ) & Условное обозначение на схеме & Значение/ Диапазон & Единица измерения & Шаг \\ 
\hline
\endfirsthead
\caption*{\hspace{3pt}\emph{Продолжение таблицы \ref{zpn:tbl1}\hfill\vspace{-0.5\baselineskip}}} \\ % сделано по ГОСТ 2.105 п.6.8.7
\hline
\rowcolor{gray!30}
Параметр (Параметр на ИЧМ) & Условное обозначение на схеме & Значение/ Диапазон & Единица измерения & Шаг \\ 
%\hline
\endhead
%\hline
\endfoot
\hline
\endlastfoot
%==+t1*PTOV1|LVTOV
\centering Ввод функции в работу (Ввод\_функции)  & \centering SGF1 & \centering Не предусмотрено \\ Предусмотрено & \centering -- & \centering \arraybackslash -- \\
\hline
\centering Напряжение срабатывания (Uср)  & \centering U> & \centering 50,0...250,0 & \centering В & \centering \arraybackslash 0,1 \\
\hline
\centering Напряжение срабатывания ОП (U2ср)  & \centering U2> & \centering 2,0...50,0 & \centering В & \centering \arraybackslash 0,1 \\
\hline
\centering Контроль напряжения ОП (Контроль\_U2)  & \centering SGF3 & \centering Не предусмотрено \\ Предусмотрено & \centering -- & \centering \arraybackslash -- \\
\hline
\centering Выбор режима работы (Реж\_работы)  & \centering SGF2 & \centering По 1 линейному \\ По 2 линейным \\ По 3 линейным & \centering -- & \centering \arraybackslash -- \\
\hline
\centering Выдержка времени срабатывания (Tср)  & \centering T1 & \centering 0...100000 & \centering мс & \centering \arraybackslash 10 \\
%===t1
\end{longtable}
\normalsize

\end{enumerate}
